% Display aliases accept both raw and escaped underscores from generated rows.
\ExplSyntaxOn
\NewDocumentCommand{\PaperCircuitLabel}{m}{
  \tl_set:Nn \l_tmpa_tl {#1}
  \tl_replace_all:Nnn \l_tmpa_tl {\_} {_}
  \str_case_e:nnF {\tl_to_str:V \l_tmpa_tl} {
    {qft_n18}{QFT-18}
    {ising_n42}{Ising-42}
    {wstate_n27}{Wstate-27}
    {ising_model_16}{Ising-16}
    {ising_model_13}{Ising-13}
    {ising_model_10}{Ising-10}
  } {#1}
}
\NewDocumentCommand{\PaperDisplayRepresentativeRows}{}{
  \tl_set:No \l_tmpb_tl {\GAMainRepresentativeCircuitRows}
  \tl_replace_all:Nnn \l_tmpb_tl {QMAP154} {QMAP}
  \tl_use:N \l_tmpb_tl
}
\ExplSyntaxOff
\begin{table}[H]
\caption{\PaperRevision{部分电路的模型保真度比较}}
\label{tab:representative-circuits}
\centering
\PaperTableSetup
\fontsize{8.5pt}{9.7pt}\selectfont
\let\PaperRawTexttt\texttt
\renewcommand{\texttt}[1]{\PaperRawTexttt{\PaperCircuitLabel{#1}}}
\begin{PaperRevisionBlock}
\begin{tabularx}{\columnwidth}{@{}lLcrrr@{}}
\toprule
基准集 & 电路 & $n/N_2$ & ZAC & ICCAD & GA-LK\\
\midrule
\PaperDisplayRepresentativeRows
\bottomrule
\end{tabularx}
\end{PaperRevisionBlock}
\PaperTableNote{\fontsize{8pt}{9.1pt}\selectfont\PaperRevision{注：$n/N_2$为量子比特数/双比特门数；ICCAD指ICCAD/QMAP，粗体为该行最优值。从两种基线较高保真度不低于0.05的电路中，每集按GA-LK相对该值的对数增幅取前两项。}}
\end{table}
